\begin{abstract}
Humans write code in a fundamentally interactive manner and rely on constant execution feedback to correct errors, resolve ambiguities, and decompose tasks. While LLMs have recently exhibited promising coding capabilities, current coding benchmarks mostly consider a static instruction-to-code sequence transduction process, which has the potential for error propagation and a disconnect between the generated code and its final execution environment. To address this gap, we introduce \benchmark{}, a lightweight, flexible, and easy-to-use framework of interactive coding as a standard reinforcement learning (RL) environment, with code as actions and execution feedback as observations. Our framework is language and platform agnostic, uses self-contained Docker environments to provide safe and reproducible execution, and is compatible out-of-the-box with traditional seq2seq coding methods, while enabling the development of new methods for interactive code generation.
We use \benchmark{} to create two interactive code environments with Bash and SQL as action spaces, leveraging data from the static Spider~\cite{yu-etal-2018-spider} and NL2Bash~\cite{lin-etal-2018-nl2bash} datasets. We demonstrate \benchmark{}'s viability as a testbed by evaluating multiple state-of-the-art LLMs configured with different prompting strategies such as ReAct~\cite{yao2023react} and Plan \& Solve~\cite{wang2023planandsolve}. Our results showcase the benefits of interactive code generation and demonstrate that \benchmark{} can serve as a challenging benchmark for advancing code understanding and generation capabilities. \benchmark{} is designed to be easily extensible and can even be used to incorporate new tasks such as Capture the Flag, a popular coding puzzle that is inherently multi-step and involves multiple programming languages.
\end{abstract}